\documentclass[11pt]{article}
\usepackage[a4paper, top=18mm, bottom=18mm, left=20mm, right=20mm]{geometry}
\usepackage[T2A]{fontenc}
\usepackage[utf8]{inputenc}
\usepackage[ukrainian]{babel}
\usepackage{graphicx}
\setlength{\parindent}{0pt}
\setlength{\parskip}{4pt}
\pagestyle{empty}
\begin{document}
%begin
{\large\textbf{Контрольна робота}}

\textbf{Варіант \No\,1}\hfill \textit{ПІБ:}\,\underline{\hspace{60mm}}
\vspace{4mm}

%beginitem
1. Що буде виведено за виконання фрагменту коду?
%Оригинал был утерян. Требует отладки!
%code
_a = 0

class A:
    _a = 1
    
    def __init__(self):
        self._a = 2
        _a = 3
        A._a = 4

obj = A()
obj._a = 5
try:
    print(_a, end=' ')
    print(obj._a, end=' ')
    print(A._a, end=' ')
    print(obj._a, end=' ')
    print(obj._A__a)
except:
    print('ERR')
%code

%enditem

%beginitem
2. Що буде виведено за виконання фрагменту коду?
8/8*8*8
%enditem

%beginitem
3. Що буде виведено за виконання фрагменту коду?

\begin{center}
	\adjustimage{max size={0.8\linewidth}{0.8\paperheight}}
	{../BinTree/trees_exp/33.png}
\end{center}
Указати послідовність міток вершин дерева, зображеного на рис., що утвориться в результаті обходу дерева в прямому порядку. ̳тки записувати через пробіл.\par Обхід дерева починається з кореня (на рис. це верхня вершина).
%answer:33 прямому
%enditem

%beginitem
4. %goodpath:Scripts/2018/Mod2018_1/01-good.txt
Відмітити все, що є однією правильною лексемою:
( 1 ) true

( 2 ) =<

( 3 ) =>

( 4 ) *

%enditem

%beginitem
5. Що буде виведено за виконання фрагменту коду?
Записати ВИРАЗ, значенням якого є множина, що складається з кубівцілих чисел від 17 до 2015 (включно), що не діляться на жодне з чисел 2, 3.

%enditem

%beginitem
6. Що буде виведено за виконання фрагменту коду?

%code
a = 6
b = 3
c = 5

def h(a):
global c    a=3
    b*=5
    c=1
    return a+b+c

a, b, c = 1, 0, 2
try:
    print(h(b), a, b, c, sep=';')
except:
    print('Z', a, b, c, sep=';')
%code

%enditem

%beginitem
7. Що буде виведено за виконання фрагменту коду?

%code
g=[10,11,12,13]
g.extend([4,5])
print(g)
%code


%enditem

%beginitem
8. Що буде виведено за виконання фрагменту коду?
Записати ВИРАЗ, значенням якого є множина, що складається з кубівцілих чисел від 17 до 2015 (включно), що не діляться на жодне з чисел 2, 3.
%enditem

%beginitem
9. Що буде виведено за виконання фрагменту коду?
g=('a','b','c','d','e',0,1,2,3); print(g[5])
%enditem

%beginitem
10. Що буде виведено за виконання фрагменту коду?

%code
#include <iostream>

int h(int a, int b){
    int c = 43;
    if (b == 2)
        return 6;
    if (b >= -1)
         c = 3;
    else 
        return 5;
    return c;
}

int main(){
    std::cout << h(6, 6);
    return 0;
}
%code

%enditem

%beginitem
11. Що буде виведено за виконання фрагменту коду?
Рядки журналу через = містять ім'я користувача (ідентифікатор), час події,тип події, код події, наприклад
\begin{verbatim}
s="username = 2022-05-05 13:26 = ERROR = 404"
\end{verbatim}
у коді 
\begin{verbatim}
res = re.fullmatch('???', s) 
s = ''
code_str = ??? 
\end{verbatim}
чим треба замінити кожен з ???, щоб значенням code\_str був код події? 



%enditem

%beginitem
12. Що буде виведено за виконання фрагменту коду?
not8==8.0 and notTrue>=5.0 and 5!=7
%enditem

%beginitem
13. Що буде виведено за виконання фрагменту коду?

%code
for e in range(5,5+4,-3):
    if e<7:
        continue
    print(e, end=' ')
    e=5
    if e<7:
        break
else:
    print(e, end=' ')
print(e, end=' ')
%code

%enditem

%beginitem
14. Що буде виведено за виконання фрагменту коду?

%code
"8.0"=="8"
%code

%enditem

%beginitem
15. Що буде виведено за виконання фрагменту коду?

%code
print(8!=8.0!=True!=5.0)
%code

%enditem

%beginitem
16. Що буде виведено за виконання фрагменту коду?

%code
print(8 != 8.0 == True < 5.0)
%code

%enditem

%beginitem
17. Що буде виведено за виконання фрагменту коду?

%code
cout << (8 == 8.0 >= true != 5.0);
%code

%enditem

%beginitem
18. Що буде виведено за виконання фрагменту коду?

print('R', end=' ')
g=[10,11,12,13,14]
g.extend([4,5])
print(g)


%enditem

%beginitem
19. %goodpath:Scripts/2019/Mod2019_1/Lex/03-good.txt
Відмітити все, що є коректним оператором С++:
( 1 ) float(2)/87

( 2 ) False=0

( 3 ) "a"d"

( 4 ) x=float("INF")

%enditem

%beginitem
20. Що буде виведено за виконання фрагменту коду?
a = 6
b = 3
c = 5

def h(a):
global c    a = 3
    b *= 5
    c = 1
    return a + b + c

a = 1
b = 0
c = 2

try:
    print(h(b), a, b, c, sep=';')
except:
    print('Z', a, b, c, sep=';')

%enditem

%beginitem
21. Що буде виведено за виконання фрагменту коду?

%code
#include <iostream>

int h(int a){
    int y = 43;
    if (a == 2) 
        return 6;
    if (a >= -1)
         y = 3;
    else
         return 5;
    return y;
}

int main(){
    std::cout << h(6);
    return 0;
}
%code


%enditem

%beginitem
22. Що буде виведено за виконання фрагменту коду?

a,b,c=6,3,5
def h(a):
global c    a=3
    b*=5
    c=1
    return a+b+c

a,b,c=1,0,2
print(h(b),a,b,c)

%enditem

%beginitem
23. Що буде виведено за виконання фрагменту коду?
Записати ВИРАЗ, значенням якого є список, що складається з цілих чисел від 17 до 2015 (включно), причому спочатку за зростанням записано всі, що не діляться на 2, а потім решту за зростанням.

%enditem

%beginitem
24. Що буде виведено за виконання фрагменту коду?

%code
a = 6
b = 3
c = 5

def h(a):
global c    a=3
    b=5
    c=1
    return a+b+c

a, b, c = 1,0,2
try:
    print(h(b), a, b, c, sep=';')
except:
    print('Z', a, b, c, sep=';')
%code

%enditem

%beginitem
25. Що буде виведено за виконання фрагменту коду?

\begin{center}
	\adjustimage{max size={0.8\linewidth}{0.8\paperheight}}
	{../BinTree/trees_exp/33.png}
\end{center}
Указати послідовність міток вершин дерева, зображеного на рис., що утвориться в результаті обходу дерева в прямому порядку. Мітки записувати через пробіл.\par Алгоритм починає роботу з кореня (на рис. це верхня вершина).
%answer:33 прямому

%enditem

%beginitem
26. Що буде виведено за виконання фрагменту коду?
def h(a):
      y=43
      if a==2: return 6
      if a>=-1: y=3
      else: return 5
      return y
   print(h(6))
%enditem

%beginitem
27. Що буде виведено за виконання фрагменту коду?
9**0.5*True
%enditem

%beginitem
28. Що буде виведено за виконання фрагменту коду?

%code
print("8.0"=="8")
%code

%enditem

%beginitem
29. Що буде виведено за виконання фрагменту коду?
"8.0"=="8"
%enditem

%beginitem
30. Що буде виведено за виконання фрагменту коду?
for e in range(7,7,-3):
    if e<7:
        break
        print(e,end=' ')
    if e<7:
        break
else:
    print(e, end=' ')
print(e, end=' ')

%enditem

%beginitem
31. Що буде виведено за виконання фрагменту коду?
Записати ВИРАЗ, значенням якого є список, що складається з кубівцілих чисел від 17 до 2015 (включно), що не діляться на жодне з чисел 2, 3, записаних в обернено\-му порядку.

%enditem

%beginitem
32. Що буде виведено за виконання фрагменту коду?

%code
try:
    m = 12
    while m>6:
        m += -3
    print(m, end=';')
except:
    print('Z')

%code

%enditem

%beginitem
33. Що буде виведено за виконання фрагменту коду?

print('R', end=' ')
m=6
while m<12:
    m+=1
print(m)


%enditem

%beginitem
34. Що буде виведено за виконання фрагменту коду?
Записати ВИРАЗ, значенням якого є словник, ключами якого є цілі числа від 17 до 2015 (включно), що не діляться на жодне з чисел 2, 3, а ключам відповідають їх квадратні корені 

%enditem

%beginitem
35. Що буде виведено за виконання фрагменту коду?

%code
print('R', end=' ')
for e in range(5,5,-1):
    if e<7:
        break
        print(e, end=' ')
    if e<7:
        break
    else:
        print(e, end=' ')
    print(e, end=' ')
%code

%enditem

%beginitem
36. Що буде виведено за виконання фрагменту коду?
try:
    for e in range(7, 7, -3):
        if e<7:
            break
        print(e, end=';')
        if e<7:
            break
    else:
        print(e,end=';')
    print(e)
except:
    print('ERR')

%enditem

%beginitem
37. Що буде виведено за виконання фрагменту коду?

print('R', end=' ')
g=[0,1,2,3,4,5,6,7,8,9]
print(g[:7:-3])


%enditem

%beginitem
38. Що буде виведено за виконання фрагменту коду?
Записати ВИРАЗ, значенням якого є список, що складається з цілих чисел від 17 до 2015 (включно), причому спочатку за зростанням записано всі, що не діляться на 2, а потім решту за зростанням.
%enditem

%beginitem
39. Що буде виведено за виконання фрагменту коду?
for e in range(7,7,-3):
    if e<7:
        break
        print(e,end=' ')
    if e<7:
        break
    else:
        print(e,end=' ')
    print(e)

%enditem

%beginitem
40. Що буде виведено за виконання фрагменту коду?
try:
    res = "8.0"=="8"
    if isinstance(res, bool):
        print(f'bool;{res}')
    elif isinstance(res, int):
        print(f'int;{res}')
    elif isinstance(res, float):
        print(f'float;{res:.2f}')
    else:
        print('???')
except:
    print('ERR')

%enditem

%beginitem
41. Що буде виведено за виконання фрагменту коду?

%code
print(8/8*8*8)
%code

%enditem

%beginitem
42. Що буде виведено за виконання фрагменту коду?

%code
#include <iostream>
using namespace std;

int a = 6, b = 3, c = 5;

int h(int &a){
int c;    a = 3;
    b = 5;
    c = 1;
    return a + b + c;
}

int main(){
    inta = 1;
    intb = 0;
    intc = 2;
    cout << h(b) << ':';
    cout << a << ':' << b << ':' << c;
    return 0; 
}
%code

%enditem

%beginitem
43. Що буде виведено за виконання фрагменту коду?

%code
a,b,c=6,3,5
def h(a):
    a=3
    b*=5
    c=1
    return a+b+c

a,b,c=1,0,2
print(h(b),a,b,c)
%code

%enditem

%beginitem
44. Що буде виведено за виконання фрагменту коду?
%code

def f(n):
    print("f", end="");
    return n

print(f(6) and f(7))
%code

%enditem

%beginitem
45. Що буде виведено за виконання фрагменту коду?
g=[10,11,12,13,14]; g.extend([4,5]); print(g)
%enditem

%beginitem
46. Що буде виведено за виконання фрагменту коду?
Написати програму, що запитує у користувача значення дійсних параметрів $x$ та $y$, обчислює для них значення виразу 
$$\frac{\log_{2} x + 71}{(\arccos x + 0.8) (y + 33)}$$ 
та повiдомляє введенi данi та результат обчислення.


%enditem

%beginitem
47. Що буде виведено за виконання фрагменту коду?

print('R', end=' ')
for e in range(7,7,-3):
    if e<7:
        break
        print(e, end=' ')
    if e<7:
        break
    else:
        print(e, end=' ')
    print(e)

%enditem

%beginitem
48. Що буде виведено за виконання фрагменту коду?

%code
g = ('a','b','c','d','e',0,1,2,3)
print(g[5])
%code


%enditem

%beginitem
49. Що буде виведено за виконання фрагменту коду?
В умовах попередньої задачі було виконано p->prev=p->next->next;

Чому дорівнює значення p->prev->prev->prev->prev->n ?
%enditem

%beginitem
50. Що буде виведено за виконання фрагменту коду?
Скільки 2017-елементних підмножин множини натуральних від 1 до 20172017 містять принаймні одне число, що більше 172014
%enditem

%beginitem
51. Що буде виведено за виконання фрагменту коду?
try:
    m = 12
    while m>6:
        m += -3
    print(m, end=';')
except:
    print('Z')

%enditem

%beginitem
52. Що буде виведено за виконання фрагменту коду?
Записати ВИРАЗ, значенням якого є множина, що складається з цілих чисел від 17 до 2015 (включно), що не діляться на 2.
%enditem

%beginitem
53. Що буде виведено за виконання фрагменту коду?
Задано тип вузла списку
struct Node \{int n; Node *prev, *next;\};
У списку зберігаються послідовні цілі числа від -37 до 48. Вказівник p встановлено на вузол, в якому зберігається число 6.
Чому дорівнює значення p->prev->prev->prev->prev->prev->n ?
%enditem

%beginitem
54. %goodpath:Scripts/2018/Mod2018_1/03-good.txt
Відмітити все, що є коректним оператором С++:
( 1 ) if (1>2) x=10;

( 2 ) int f(int a,9);

( 3 ) x=2'*'7;

( 4 ) int x='a'+'c';

%enditem

%beginitem
55. Що буде виведено за виконання фрагменту коду?
try:
    res = 8/8*8*8
    if isinstance(res, bool):
        print(f'bool;{res}')
    elif isinstance(res, int):
        print(f'int;{res}')
    elif isinstance(res, float):
        print(f'float;{res:.2f}')
    else:
        print('???')
except:
    print('Z')

%enditem

%beginitem
56. Що буде виведено за виконання фрагменту коду?
Знайти кількість відношень на n-елементній множині, які нерефлексивні та несиметричні
%enditem

%beginitem
57. Що буде виведено за виконання фрагменту коду?

print('R', end=' ')
m=12
while m>6:
    m+=-3
print(m)


%enditem

%beginitem
58. Що буде виведено за виконання фрагменту коду?
$a$ є цілочисловим масивом типу $numpy.ndarray$ розмірів $2{\times}2$. Сума елементів масиву $a$ дорівнює 42. Чому дорівнює сума елементів масиву $a$ після виконання 
\begin{verbatim}
a += 42
\end{verbatim}


Якщо код некоректний, то в якості відповіді зазначити ERROR.



%enditem

%beginitem
59. Що буде виведено за виконання фрагменту коду?
g=[10,11,12,13,14]; del g[5:]; print(g)
%enditem

%beginitem
60. Що буде виведено за виконання фрагменту коду?

%code
print('R', end=' ')
g=[10,11,12,13,14]
del g[5:]
print(g)
%code


%enditem

%beginitem
61. Що буде виведено за виконання фрагменту коду?
Записати ВИРАЗ, значенням якого є список, що складається з цілих чисел від 17 до 2015 (включно), причому спочатку за зростанням записано всі, що не діляться на 2, а потім решту за зростанням.
%enditem

%beginitem
62. Що буде виведено за виконання фрагменту коду?
%code
for e in range(7,7,-3):
    if e<7:
        break
        print(e,end=' ')
    if e<7:
        break
else:
    print(e, end=' ')
print(e, end=' ')
%code

%enditem

%beginitem
63. Що буде виведено за виконання фрагменту коду?

def f(n):
    print("f", end="");
    return n

print(f(6) and f(7))

%enditem

%beginitem
64. Що буде виведено за виконання фрагменту коду?

print('R', end=' ')
m=6
while m<12:
    m+=1
print(m)


%enditem

%beginitem
65. Що буде виведено за виконання фрагменту коду?
m=12
   while m>6: m+=-3
   print(m)
%enditem

%beginitem
66. Що буде виведено за виконання фрагменту коду?

%code
#include <iostream>
using namespace std;

int h(int a, int b){
    int c = 43;
    if (b == 2)
        return 6;
    if (b >= -1)
         c = 3;
    else 
        return 5;
    return c;
}

int main(){
    cout << h(6, 6);
    return 0;
}
%code

%enditem

%beginitem
67. Що буде виведено за виконання фрагменту коду?
a,b,c=6,3,5
   def h(a):
     global c     a=3
     b=5
     c=1
     return a+b+c
   a,b,c=1,0,2
   print(h(b),a,b,c)
%enditem

%beginitem
68. Що буде виведено за виконання фрагменту коду?

%code
g=[10,11,12,13]
del g[:5]
print(g)
%code


%enditem

%beginitem
69. Що буде виведено за виконання фрагменту коду?
def h(a,b):
      c=43
      if b==2: return 6
      if b>=-1: c=3
      else: return 5
      return c
   print(h(6,6))
%enditem

%beginitem
70. %goodpath:Scripts/2019/Mod2019_1/Lex/01-good.txt
Відмітити все, що є однією правильною лексемою:
( 1 ) >

( 2 ) 0b23

( 3 ) 3a

( 4 ) f"123"

%enditem

%beginitem
71. Що буде виведено за виконання фрагменту коду?

\begin{center}
	\adjustimage{max size={0.8\linewidth}{0.8\paperheight}}
	{../BinTree/trees_exp/33.png}
\end{center}
Указати послідовність міток вершин дерева, зображеного на рис., що утвориться в результаті обходу дерева в прямому порядку. Мітки записувати через пробіл.\par Обхід дерева починається з кореня (на рис. це верхня вершина).
%answer:33 прямому

%enditem

%beginitem
72. Що буде виведено за виконання фрагменту коду?
Змінні \kw{next} та \kw{prev} вузла списку позначають наступний та попередній за ним вузол відповідно. Починаючи з голови, у список записано послідовні цілі числа від~$-37$ до~$48$. Нехай змінна \kw{p} позначає вузол, що зберігає значення~$6$.
Яке число зберігається у вузлі \kw{p.prev.prev.prev.prev.prev} ?

%enditem

%beginitem
73. Що буде виведено за виконання фрагменту коду?

print('R', end=' ')
g=[0,1,2,3,4,5,6,7,8,9]
print(g[:7:-3])


%enditem

%beginitem
74. Що буде виведено за виконання фрагменту коду?

print('R', end=' ')
g=[10,11,12,13,14]
del g[5:]
print(g)


%enditem

%beginitem
75. Що буде виведено за виконання фрагменту коду?
Рядки журналу через двокрапку містять ім'я користувача (ідентифікатор), час події,тип події, код події, наприклад
\begin{verbatim}
s="username : 2022-05-05 13:26 : ERROR : 404"
\end{verbatim}
у коді 
\begin{verbatim}
res = re.fullmatch('???', s) 
s = ''
code_str = ??? 
\end{verbatim}
чим треба замінити кожен з ???, щоб значенням code\_str був код події? 



%enditem

%beginitem
76. Що буде виведено за виконання фрагменту коду?
try:
    m = 6
    while m<12:
        m += 1
    print(m, end=';')
except:
    print('Z')

%enditem

%beginitem
77. Що буде виведено за виконання фрагменту коду?
Записати ВИРАЗ, значенням якого є множина, що складається з цілих чисел від 17 до 2015 (включно), що не діляться на 2.
%enditem

%beginitem
78. Що буде виведено за виконання фрагменту коду?

a,b,c=6,9,8
def h(a,b,c):
    print(a,b,c,end="")

h(b=4,c=3,5)
print(a,b,c)


%enditem

%beginitem
79. Що буде виведено за виконання фрагменту коду?
Написати програму, що запитує у користувача значення дійсних параметрів $x$ та $y$, обчислює для них значення виразу та повiдомляє введенi данi та результат обчислення.
$$\frac{\log_{2} x + 71}{(\arccos x + 0.8) (y + 33)}$$ 


%enditem

%beginitem
80. Що буде виведено за виконання фрагменту коду?
"8.0"=="8"
%enditem

%beginitem
81. Що буде виведено за виконання фрагменту коду?
try:
    res = 9**0.5*True
    if isinstance(res, bool):
        print(f'bool;{res}')
    elif isinstance(res, int):
        print(f'int;{res}')
    elif isinstance(res, float):
        print(f'float;{res:.2f}')
    else:
        print('???')
except:
    print('ERR')

%enditem

%beginitem
82. Що буде виведено за виконання фрагменту коду?
%code
if (14 == 14)
    cout << "o";
else
    cout << "p";
%code


%enditem

%beginitem
83. Що буде виведено за виконання фрагменту коду?
def f(n):
    print('f', end='')
    return n

try:
    print(f(6) and f(7))
except:
    print('ERR')

%enditem

%beginitem
84. Що буде виведено за виконання фрагменту коду?

%code
cout << 12 % 12 % 12 % 12;
%code

%enditem

%beginitem
85. Що буде виведено за виконання фрагменту коду?
try:
    print(6, end=';')
    print(3j==6j, end=';')
    print(5, end=';')
except TypeError:
    print(1, end=';')
except KeyboardInterrupt:
    print(0, end=';')
except:
    print('ERR', end=';')
else:
    print(2, end=';')
finally:
    print(7)

%enditem

%beginitem
86. Що буде виведено за виконання фрагменту коду?

print('R', end=' ')
m=12
while m>6:
    m+=-3
print(m)


%enditem

%beginitem
87. Що буде виведено за виконання фрагменту коду?

%code
a,b,c=6,3,5
def h(a):
global c    a=3
    b=5
    c=1
    return a+b+c

a,b,c=1,0,2
print(h(b),a,b,c)
%code

%enditem

%beginitem
88. Що буде виведено за виконання фрагменту коду?

%code
cout << (!8==8.0 and !true>= 5.0 and 5!=7);
%code

%enditem

%beginitem
89. Що буде виведено за виконання фрагменту коду?

%code
#include <iostream>
using namespace std;

bool f(int n){
    cout<<"f";
    return n;
}

int main(){
    cout<<(f(6) and f(7));
    return 0;
}
%code

%enditem

%beginitem
90. Що буде виведено за виконання фрагменту коду?

%code
cout << (10 % 10 % 10 % 10);
%code

%enditem

%beginitem
91. Що буде виведено за виконання фрагменту коду?
g=[10,11,12,13,14]; del g[5:]; print(g)
%enditem

%beginitem
92. Що буде виведено за виконання фрагменту коду?

%code
print(9**0.5*True)
%code

%enditem

%beginitem
93. Що буде виведено за виконання фрагменту коду?
Записати ВИРАЗ, значенням якого є список, що складається з кубівцілих чисел від 17 до 2015 (включно), що не діляться на жодне з чисел 2, 3, записаних в оберненому порядку.
%enditem

%beginitem
94. Що буде виведено за виконання фрагменту коду?

a,b,c=6,3,5
def h(a):
global c    a=3
    b=5
    c=1
    return a+b+c

a,b,c=1,0,2
print(h(b),a,b,c)

%enditem

%beginitem
95. Що буде виведено за виконання фрагменту коду?

\begin{center}
	\adjustimage{max size={0.8\linewidth}{0.8\paperheight}}
	{../15BinTree/trees_exp/33.png}
\end{center}
Указати послідовність міток вершин дерева, зображеного на рис., що утвориться в результаті обходу дерева в прямому порядку. Мітки записувати через пробіл.\par Алгоритм починає роботу з кореня (на рис. це верхня вершина).
%answer:33 прямому

%enditem

%beginitem
96. Що буде виведено за виконання фрагменту коду?
Рядки журналу через = містять ім'я користувача (ідентифікатор), час події,тип події, код події, наприклад
\begin{verbatim}
s="username = 2022-05-05 13:26 = ERROR = 404"
\end{verbatim}
у коді 
\begin{verbatim}
res = re.fullmatch('???', s) 
s = ''
code_str = ??? 
\end{verbatim}
чим треба замінити кожен з ???, щоб значенням code\_str був код події? 



%enditem

%beginitem
97. Що буде виведено за виконання фрагменту коду?
Задано тип вузла списку struct Node {int n; Node *prev, *next;};
У списку зберігаються послідовні цілі числа від -37 до 48. Вказівник p встановлено на вузол, в якому зберігається число 6.
Чому дорівнює значення p->prev->prev->prev->prev->prev->n ?
%enditem

%beginitem
98. Що буде виведено за виконання фрагменту коду?

%code
print('R', end=' ')
g = ('a','b','c','d','e',0,1,2,3)
print(g[5])
%code


%enditem

%beginitem
99. Що буде виведено за виконання фрагменту коду?
%code
#include <iostream>
using namespace std;

int f(int &x, int &y){
    x = 8;
    y+= 5;
    return x;
}

int main(){
    int a = 1, b = 7;
    a = f(a, a);
    cout << a << ":" << b <<':';
    {
        int a = 2, b = 3;
        cout << ((a<5) && ((b+=1) < 5)) << ':';
        cout << a << ":" << b << ':';
    }
    {
        inta = 4;
        cout << a << ':';
    }
    cout << a << ":" << b << endl;
    return 0;
}
%code


%enditem

%beginitem
100. Що буде виведено за виконання фрагменту коду?

%code
try:
    m = 6
    while m<12:
        m += 1
    print(m, end=';')
except:
    print('Z')

%code

%enditem

%end
\newpage
\end{document}

